\documentclass{article}
\usepackage{graphicx} % Required for inserting images

\title{Quantum Mechanics}
\author{Liam McGrath}
\date{November 2023}

\begin{document}

\maketitle

\section{Quantum Mechanics Solutions}

\subsection{Hydrogen Atom}

The Hydrogen Atom has the following Hamiltonian;

\begin{equation}
H = -\frac{\hbar^2}{2m}\nabla^2 - \frac{e^2}{4\pi\epsilon_0r}
\end{equation}

where m is the mass of the electron, e is the charge of the electron, r is the distance between the electron and the proton, and \begin{math} \epsilon_0 \end{math} is the permittivity of free space. We stand to benefit from using the spherical coordinates of the Laplacian and then separation of variables on the radial and angular components:

\begin{equation}
H = -\frac{\hbar^2}{2m}\left[\frac{1}{r^2}\frac{\partial}{\partial r}\left(r^2\frac{\partial}{\partial r}\right) + \frac{1}{r^2\sin\theta}\frac{\partial}{\partial\theta}\left(\sin\theta\frac{\partial}{\partial\theta}\right) + \frac{1}{r^2\sin^2\theta}\frac{\partial^2}{\partial\phi^2}\right] - \frac{e^2}{4\pi\epsilon_0r}
\end{equation}

Using separation of variables; \begin{math} \psi(r,\theta,\phi) = R(r)Y(\theta,\phi) \end{math} then substituting:

\begin{equation}
\frac{1}{R}\frac{d^2R}{dr^2} + \frac{2}{r}\frac{1}{R}\frac{dR}{dr} - \frac{l(l+1)}{r^2} + \frac{2mr^2}{\hbar^2}\left[E + \frac{e^2}{4\pi\epsilon_0r} - \frac{\hbar^2}{2m}\frac{l(l+1)}{r^2}\right] = 0
\end{equation}

...we get equations for the radial and angular parts of the wavefunction:

\begin{equation}
\frac{1}{Y}\left[\frac{1}{\sin\theta}\frac{\partial}{\partial\theta}\left(\sin\theta\frac{\partial Y}{\partial\theta}\right) + \frac{1}{\sin^2\theta}\frac{\partial^2 Y}{\partial\phi^2}\right] = -\frac{l(l+1)}{\hbar^2}Y
\end{equation}

where l is the angular momentum quantum number. Spherical Harmonics help us find the angular part of the solution, where m is the magnetic quantum number and \begin{math} P_l^m \end{math} is the associated Legendre polynomial.

\begin{equation}
Y_{lm}(\theta,\phi) = (-1)^m\sqrt{\frac{(2l+1)(l-m)!}{4\pi(l+m)!}}P_l^m(\cos\theta)e^{im\phi}
\end{equation}

\vspace{0.5cm}

We can solve the radial part of the Schrödinger equation by using the radial wave function:

\begin{equation}
R_{nl}(r) = \frac{1}{r}\rho^{l+1}e^{-\rho}F(-n+l+1,2l+2,\rho)
\end{equation}

...where n is the principal quantum number, l is the angular momentum quantum number, F is the hypergeometric function, and \begin{math} \rho \end{math} is a dimensionless variable defined as \begin{math} \rho = 2r/na_0 \end{math} , where \begin{math} a_0 \end{math} is the Bohr radius. The energy eigenvalues of the hydrogen atom are given by;

\begin{equation} 
E_n = -\frac{me^4}{32\pi^2\epsilon_0^2\hbar^2n^2} 
\end{equation}

and the final wavefunction is of the form, where n is the principal quantum number;

\begin{equation}
 \psi(r,\theta,\phi) = R(r)Y(\theta,\phi)
\end{equation}

\begin{equation}
 \psi(r,\theta,\phi)  = \frac{1}{r}\rho^{l+1}e^{-\rho}F(-n+l+1,2l+2,\rho)(-1)^m\sqrt{\frac{(2l+1)(l-m)!}{4\pi(l+m)!}}P_l^m(\cos\theta)e^{im\phi}
\end{equation}

\end{document}